\chapter{Integration as a Checker}\label{ch:integrate}

Nobody proves an integrator correct. The engine does not try. What it does instead is the oldest
trick in the subject --- differentiate the answer and see --- made into a theorem, and the
interesting part is what had to change in the pipeline for the theorem to be stateable at all.

\section{A guess, then a check}

The finder (\file{Antiderivative.lean}) is an ordinary table-driven antidifferentiator: sums,
constant factors, powers, $a^u$, the elementary functions, each under a linear substitution;
later, $u$-substitution and integration by parts with the LIATE heuristic. Nothing about it is
proved, and its steps are labelled with a fourth status, \status{checked}: a guess whose result a
later step verifies.

The command then differentiates the candidate with the pipeline, and accepts only if the normal
form of the derivative is the normal form of the integrand, exactly. The claim is a theorem about
the command's output, and it assumes nothing about the finder:

\begin{minted}{lean}
theorem cmdIntegrate_spec (norm : Norm) {f : Expr} {x : String} {res : RuleResult}
    (h : (cmdIntegrate norm).apply (.fn "integrate" [f, .var x]) = some res)
    (hok : res.error = none) :
    ∃ g sub g' s₁ s₂, norm (D res.result x) = .ok (g, sub) ∧
      norm (Expand.dist g) = .ok (g', s₁) ∧ norm (Expand.dist f) = .ok (g', s₂)
\end{minted}

And in \file{proofs/}, with the real semantics, that becomes a statement about \lc{deriv}:

\begin{minted}{lean}
theorem integrate_deriv (norm : Norm) … (ρ : EnvR)
    (hdiff  : … evalD ρ (D res.result x) = evalD ρ g)
    (hleft  : … evalD ρ (Expand.dist g) = evalD ρ g')
    (hright : … evalD ρ (Expand.dist f) = evalD ρ f') :
    deriv (fx ρ x res.result) (ρ x) = evalD ρ f
\end{minted}

The three hypotheses say that the normalizations the check performed were sound at the point
$\rho$. That is exactly what the statuses of the nested derivation say, rule by rule
(Chapters~\ref{ch:sound} and~\ref{ch:deriv}). So the notebook's display of an integration ---
finder steps marked \status{checked}, an \rn{int.check} step with the differentiation nested
under it, and the command inheriting the weakest status below it --- is a faithful rendering of
this theorem's hypotheses.

\section{The knot}

The pipeline could not check with itself. A rule set cannot contain a rule that normalizes with
that set --- the definition would be circular, and Lean says so. The way out was to make the
pipeline \emph{generic in the checker's normalizer}:

\begin{minted}{lean}
def pipelineRulesWith (norm : Norm) : List PlainRule := …
theorem pipelineOrderedWith (norm : Norm) : Ordered (pipelineRulesWith norm)
\end{minted}

The termination theorem of Chapter~\ref{ch:order} then holds for \emph{every} \lc{norm}, and the
knot closes after the proof: the checker is the pipeline with nested \texttt{integrate} refused,
the notebook's pipeline is the pipeline with that checker, and neither has a step budget.

\begin{minted}{lean}
def pipelineRules : List PlainRule := pipelineRulesWith checkNorm
theorem pipelineOrdered : Ordered pipelineRules := pipelineOrderedWith checkNorm
\end{minted}

\section{Exact means refusable}

The check is exact, so a correct guess can be refused. $\int \tan x$ was: the finder produced
$-\ln \cos x$, its derivative normalized to $\sin x \cdot (\cos x)^{-1}$, and the engine did not
identify that with $\tan x$. That is the right failure mode --- never a wrong answer --- and the fix
was a rule (Chapter~\ref{ch:cannot}), proved in three layers, not a weaker check.

Integration by parts found a subtler one. Every candidate was correct and every one was refused,
because the pipeline never distributes a numeral over a sum: the ordering forbids it, distribution
is \texttt{expand}'s job, and so $-(a + b) + b$ is a normal form. The fix keeps the check sound
rather than weakening it: both sides are expanded with \lc{Expand.dist} --- proved sound for the
derivative semantics too, \lc{dist_soundD}, the second instance of the generic proof --- and
simplified before comparison. That is the \rn{int.compare} step, and it is why the spec above
compares \lc{norm (Expand.dist g)} with \lc{norm (Expand.dist f)} rather than \lc{g} with
\lc{f}.

A third finding came free. $\int x^a$ is accepted through \rn{simp.collect-powers} cancelling
$(a+1)/(a+1)$, so its check inherits that rule's side condition, $a \neq -1$, and the statuses say
so without anyone writing it down. The theorem's hypotheses are the notebook's display; the
display knew about $a = -1$ before the author did.

\begin{logbox}[Still refused]
  $e^x \sin x$ needs the ``solve for the integral'' trick; $\sin^2 x$ needs a trigonometric
  identity the simplifier lacks; $1/(x^2 + 1)$ needs an arctangent the engine does not have. Each
  is refused with the candidate it tried, and none is a wrong answer.
\end{logbox}

\begin{trybox}
\begin{minted}{text}
integrate(x * exp(x), x)
integrate(x^a, x)
integrate(sin(x)^2, x)
\end{minted}
  Open the derivation of the first: the by-parts steps are \status{checked}, the check step
  carries the differentiation, and the compare step shows both sides expanded. The second's
  status names the side condition. The third is refused.
\end{trybox}

\section{Exercises}

\begin{exercisebox}[Reading the hypotheses]
  For $\int x^a\,dx$, list the rules that fire in the check's differentiation and say which
  hypothesis of \lc{integrate_deriv} each one's status discharges.
\end{exercisebox}

\begin{exercisebox}[A checker for something else]
  Design a checked command for solving $ax + b = 0$: what does the finder guess, what does the
  check normalize, and what is the analogue of \lc{cmdIntegrate_spec}?
\end{exercisebox}
